\documentclass[11pt,a4paper]{article}

\usepackage[T1]{fontenc}
\usepackage[utf8]{inputenc}
\usepackage{lmodern}
\usepackage[margin=25mm]{geometry}
\usepackage{microtype}
\usepackage{xcolor}
\usepackage{enumitem}
\usepackage{fancyvrb}
\usepackage[hidelinks]{hyperref}

\definecolor{commentbar}{gray}{0.55}
\newenvironment{reviewercomment}
  {\par\smallskip\begin{list}{}{\leftmargin=1.2em\rightmargin=0pt}\item[]\color{black!75}\itshape}
  {\end{list}\smallskip}
\newcommand{\response}{\par\noindent\textbf{Response.}~}
\newcommand{\where}[1]{\par\noindent\textbf{Revised manuscript:}~#1\par\medskip}
\setlength{\parindent}{0pt}
\setlength{\parskip}{0.4em}

\title{Response to the Reviewers\\[0.3em]
  \large gendertagger: A Multi-Method Workflow for Gender Annotation in Digital Scholarly Editions}
\author{Lou Klappenbach, Frederike Neuber, Tim Westphal}
\date{}

\begin{document}
\maketitle

We thank the three reviewers for their careful reading and constructive comments. This document
responds to each comment in turn. Reviewer comments are quoted in italics. Page and line numbers
refer to the revised manuscript unless stated otherwise; footnotes and tables carry no line numbers and
are referenced by page. Page and line numbers inside quoted reviewer comments refer to the submitted
version.

The main changes are:
\begin{itemize}[nosep]
  \item a new Discussion section (Section~5) that connects the case study with the limitations from
        a research software engineering perspective, including the role of TS4NFDI;
  \item an explanation of how the TEI \texttt{@cert} values are derived, what they mean, and why we
        chose this scale (Sections~2.2 and~2.5, Table~5);
  \item a justification of the GND as the sole authority source and a discussion of alternatives;
  \item English translations for all German examples and labels;
  \item an update of the software repository, in progress, with QUICKSTART instructions and sample
        data (see ``Software release'').
\end{itemize}

% ======================================================================
\section*{Software release (Reviewers C and K)}

Reviewers C and K both reported that the application could not be started from the repository. We
answer these comments together.

\textbf{Reviewer C:}
\begin{reviewercomment}
I tried to run the steamlit app in the repository but got an error:
\begin{Verbatim}[fontsize=\small]
from gender_project.classification.gender import (
...
ModuleNotFoundError: No module named 'gender_project'
\end{Verbatim}
Reading the README and QUICKSTART instructions, it seems like there is some information missing as
well to be able to successfully run the app; e.g.\ the instruction say to import data and provide a
default path, but that path does not seem to exist in the repository
(\path{data/raw/schleiermacher/Register/Personen}), so this would probably not work.
\end{reviewercomment}
\begin{reviewercomment}
My biggest concern is that the code can't be run when following the instructions in the README.
This should be fixed before publication.
\end{reviewercomment}

\textbf{Reviewer K:}
\begin{reviewercomment}
Unfortunately the raw data for the person register itself seems to be missing, so the streamlit
prototype has nothing to run on. (To start the prototype imports in
\path{gendertagger/src/gendertagger/prototype/app.py} have to be fixed)
\end{reviewercomment}
\begin{reviewercomment}
Unfortunately the streamlit app to showcase the workflow doesn't work without fixing the broken
imports. Even after that the workflow can't be started as described, as the needed person register
raw data is missing. If its of sensible nature, maybe you can provide a exemplary dummy file? Before
publication the instruction in the readme.md should be reproducible.
\end{reviewercomment}

\response We apologize for the delay in providing a runnable version. The Streamlit prototype
referenced in the submitted version has since been replaced: the pipeline has been reimplemented as a
local service with a FastAPI backend and a Vue frontend, and Docker containers bundle the components.
We are updating the repository for publication in the coming days.
The updated repository will include QUICKSTART instructions for installing and running the
application, and sample data.

% ======================================================================
\section*{Reviewer C}

Reviewer C's comments on running the application are answered under ``Software release'' above.

\subsection*{C.1 Software citation}

\begin{reviewercomment}
I would also recommend adding a CFF file to the repository and then include the repository in the
bibliography.
\end{reviewercomment}

\response We added a \texttt{CITATION.cff} file with the software citation metadata and included the
gendertagger repository in the bibliography. The Data Availability statement, the Reproducibility
subsection, and the Conclusion now cite this entry.

\where{p.~19, ll.~573--574; p.~21, l.~625; p.~21, ll.~647--648; bibliography entry [KNW26].}

\subsection*{C.2 Annotated PDF}

\begin{reviewercomment}
(Page 6) You might want to add a link here or short explanation what this is the abbreviation for as
I'm not sure it's well-known outside of the domain or in other countries. --- I see now that it's
described later on, but maybe just give the full name here?
\end{reviewercomment}

\response We now expand GND as the German Integrated Authority File (\emph{Gemeinsame Normdatei}) at
its first occurrence in the Introduction, together with a short explanation of authority files.

\where{p.~1, ll.~32--34.}

\begin{reviewercomment}
(Page 7) Might be worth mentioning the flip side as well; that this approach means the dump needs to
be updated every so often to include updates?
\end{reviewercomment}

\response We added this trade-off: the local snapshot must be refreshed to include later additions
and corrections to the GND. We also describe how the current backend keeps traceability across
updates: it retains source dumps and derived lookup snapshots with checksums and compilation
provenance, and records the exact lookup used by each method execution.

\where{p.~8, ll.~267--272.}

\begin{reviewercomment}
(Page 7) Is there separate code for this or is this part of gendertagger?
\end{reviewercomment}

\response The parenthetical-qualifier heuristic is implemented within gendertagger; the revised text
states this explicitly.

\where{p.~8, l.~278.}

\begin{reviewercomment}
(Page 10) Not sure I understand what the ``combined'' label is. --- Ah, after reading I think I know
what is referred to here. Maybe ``post-voting'' label or ``resulting'' label might be more accurate? ``combine'' sounds to me
like two things are put together to form a new thing, rather than one thing is chosen over another.
\end{reviewercomment}

\response We replaced ``combined label'' with ``resulting label'' throughout and define it at first use
as the automated
output of aggregation, including any authority override, before human review. Table headings, the
row label, and captions were updated accordingly.

\where{p.~7, l.~242; p.~17, l.~534; Tables~2, 6, and~7.}

% ======================================================================
\section*{Reviewer I}

\subsection*{I.0 General assessment}

\begin{reviewercomment}
Its broader impact is, however, somewhat constrained by its concentration on the GND alone, rather
than also drawing on other terminologies and community hubs such as Wikidata or international
authorities like VIAF or the Getty ULAN, which would widen its reach. [\dots] the choice of a single
authority file, without an alternative, would benefit from an explicit justification or from
consideration of other terminologies.
\end{reviewercomment}

\response See our responses to I2 and I14: the revised manuscript justifies the GND by the
predominantly German-language sources and the existing editorial GND links, explains why VIAF was
excluded, and names Wikidata and Getty ULAN as possible extensions.

\begin{reviewercomment}
In particular, a dedicated discussion of the points and decisions made in Sections 3 and 4 would
considerably strengthen the paper, especially from an RSE perspective. A recurring issue is that
numerous German phrases in the examples and tables are left untranslated [\dots]
\end{reviewercomment}

\response See I23 for the new Discussion section. Beyond the locations listed in I10 and I17--I22,
we checked all remaining German examples, the TEI example, all tables, and both figures, and added
English translations alongside the German originals. Where grammatical gender or a missing space
explains a classification error, the translation makes the distinction explicit.

\where{among others p.~6, ll.~188--191; p.~16, ll.~515--521; Tables~6 and~7.}

\subsection*{Specific comments}

\begin{enumerate}[label=\textbf{I\arabic*}, leftmargin=2.2em, itemsep=0.8em]

\item
\begin{reviewercomment}
Abstract: are the separate paragraphs needed? What does N = 8,277 refer to, and is it important to
state here?
\end{reviewercomment}
\response We merged the abstract into a single paragraph and replaced ``$N = 8{,}277$'' with
``comprising 8,277 entries'', which states that the number counts the entries of the case-study
person register. We kept the number because it indicates the scale of the study.
\where{p.~1, l.~18.}

\item
\begin{reviewercomment}
Line 32: ``Authority files, for instance, [\dots]'' -- what is meant by ``authority files'' here?
Only the GND (Gemeinsame Normdatei)? It would be preferable to use ``terminologies'' in line with NFDI
wording, which also covers thesauri, Wikidata and so on. If only the GND is meant, name it explicitly
and explain why other terminologies are not used.
\end{reviewercomment}
\response We now define authority files and name the GND as the authority source used in this study.
We kept the more specific term ``authority files'' because the method retrieves records about
identified entities, whereas ``terminologies'' covers a broader range of resources. The methods
section explains the choice of the GND and discusses VIAF, Wikidata, and Getty ULAN (see I14). We
also clarify that gender information is optional in GND person records, citing the official field
documentation.
\where{p.~1, ll.~32--36.}

\item
\begin{reviewercomment}
Footnote 1: currently runs onto the next page; it should fit on a single page.
\end{reviewercomment}
\response We changed the LaTeX footnote setting so that footnotes no longer split across pages.
Footnote~1 now appears complete on a single page.
\where{p.~2, footnote~1.}

\item
\begin{reviewercomment}
Line 77: [\dots] a link to TEI in the footnote would be very helpful.
\end{reviewercomment}
\response We added a footnote with the link to the Text Encoding Initiative website and an access date.
\where{p.~3, footnote~2.}

\item
\begin{reviewercomment}
Line 79: ``[\dots] multiple classification approaches: authority data queries, heuristic methods,
[\dots]'' -- what are ``authority data queries''?
\end{reviewercomment}
\response We replaced the phrase with an explicit description: retrieval of recorded gender
information from a local GND snapshot using the register's GND identifiers.
\where{p.~3, ll.~83--84.}

\item
\begin{reviewercomment}
Line 80: ``[\dots] Transformer-based models, [\dots]'' -- should ``Transformer'' be lower-case?
\end{reviewercomment}
\response Yes. We changed ``transformer-based'' to lower case throughout the running text and kept the
capitalization in the section heading.
\where{p.~1, l.~14; p.~3, l.~85; p.~7, l.~252; p.~20, l.~613.}

\item
\begin{reviewercomment}
Line 83: why is ``Schleiermacher digital'' the case study here? What makes it particularly suitable,
and why not another?
\end{reviewercomment}
\response We added the rationale. The edition combines a manageable corpus with enrichment challenges
that are relevant to our planned work on other Academy editions: its temporal and regional focus is
comparable, and incomplete names, relational descriptions, and limited authority coverage test
complementary methods. The collaboration with its editorial team also addressed a practical need as
the project approaches the end of its term.
\where{p.~3, ll.~88--95.}

\item
\begin{reviewercomment}
Lines 91--97: consolidate into a single paragraph.
\end{reviewercomment}
\response We merged the opening passages of the Study Design section into one paragraph.
\where{p.~3, ll.~103--109.}

\item
\begin{reviewercomment}
Lines 101--128: instead of ``First, \dots'', ``Second, \dots'', please use LaTeX
\texttt{\textbackslash paragraph\{\dots\}}.
\end{reviewercomment}
\response We replaced the ordinal openings with five descriptive paragraph headings, one for each
design constraint.
\where{p.~4, ll.~113--141.}

\item
\begin{reviewercomment}
Line 111: ``uneindeutig'' -- all other tags here are in English, so an English translation would help.
\end{reviewercomment}
\response We added the translation ``[ambiguous]'' and kept the original label, since it is the raw
output of the method.
\where{p.~4, l.~124.}

\item
\begin{reviewercomment}
Footnote 2: please add ``(accessed on DATE)''.
\end{reviewercomment}
\response Added.
\where{p.~3, footnote~3.}

\item
\begin{reviewercomment}
Lines 167--169: from a geoinformatics perspective, the snippet
``\texttt{<note>Lehrerin in <placeName key="S0000357">Gnadenfrei</placeName></note>}'' could make use of
the cert="\dots" attribute; a link to a gazetteer would also be helpful.
\end{reviewercomment}
\response The note and its place reference are reproduced unchanged from the schleiermacher digital
register; gendertagger adds only the \texttt{<gender>} element, and its \texttt{@cert} qualifies the
gender annotation alone. We did not add \texttt{@cert} to the place reference: identifying places,
and stating the certainty of that identification, remain editorial decisions of the edition, and
changing the source markup would misrepresent the pipeline's output. The edition already links this
place to a gazetteer: the key \texttt{S0000357} identifies its register entry for Gnadenfrei, which
references GeoNames record 3089025. The revised text states which markup gendertagger adds and gives
both links in a footnote.
\where{p.~5, ll.~184--187, and footnote~5.}

\item
\begin{reviewercomment}
Lines 186 ff.: how are the @cert categories derived? What do high / low or unknown mean? Would another
scale be possible, and which scales are used elsewhere in the DH domain?
\end{reviewercomment}
\response \texttt{@cert} states the certainty of the resulting annotation, not the confidence of a
single classifier. It is derived from the pipeline's certainty tier (Section~2.5): \texttt{high} when
at least two methods vote and at least three quarters agree, when an authority signal (GND or the
parenthetical-qualifier rule) decides the label, or after human review; \texttt{low} when at least
half but fewer than three quarters agree; \texttt{unknown} when the methods reach no majority or only
one method provides a label. Following the TEI definition, \texttt{unknown} means that the encoder does
not wish to assert an opinion. We also corrected the definition of the tiers: an authority override always yields
\texttt{high} (the submitted version said ``in agreement with the vote''), an unresolved tie between
substantive labels is assigned to \texttt{contested}, and the \texttt{@cert} description now covers
entries with a single voting method.

Other scales are possible: TEI also allows \texttt{medium} and numerical probabilities. We did not use
numerical values because the pipeline has no calibrated probability: per-method confidences are not
comparable, and the agreement score measures consensus, not correctness. We did not use
\texttt{medium} because the pipeline distinguishes only two levels of agreement; this coarse scale
also matches editorial practice such as the ediarum framework, which offers only \texttt{low} and
\texttt{high} for uncertain readings. The finer tiers and agreement scores remain available in the pipeline output, and
Table~5 now shows the TEI value for each tier.

While revising, we corrected passages that had called internal tiers TEI values, and the certainty
tiers given in Table~6 (see ``Further changes'').
\where{p.~6, ll.~188--217; p.~11, ll.~389--398; p.~15, l.~488; Tables~5 and~6; p.~19, l.~566.}

\item
\begin{reviewercomment}
Lines 238--241: [\dots] why only the GND? What about Wikidata or other international person
terminologies?
\end{reviewercomment}
\response We explain that the choice of the GND reflects the predominantly German-language source
material and the existing editorial authority links, which primarily use GND identifiers. VIAF was
excluded because of ongoing discussion about reliability and service governance following its January
2025 redesign (Bargioni, 2025). We name the use of records in Wikidata and Getty ULAN as possible
future extensions.
\where{p.~8, ll.~272--277.}

\item
\begin{reviewercomment}
Footnotes 5/6: add ``(accessed on DATE)''.
\end{reviewercomment}
\response Added.
\where{p.~9 and p.~10, footnotes.}

\item
\begin{reviewercomment}
Line 306: place the link as a footnote rather than in the running text.
\end{reviewercomment}
\response We moved the internal reference to the TEI example into a footnote.
\where{p.~10, footnote.}

\item
\begin{reviewercomment}
Line 347: ``[\dots] inserted by the human editors (e.g.\ Tochter des \dots), and [\dots]'' --
``Tochter des'' needs a translation, perhaps as a footnote.
\end{reviewercomment}
\response We added the translation in square brackets directly after the German phrase; the same
convention is used for all German examples.
\where{p.~11, ll.~383--384.}

\item
\begin{reviewercomment}
Line 400: provide a translation for ``[\dots] such as e.g.\ Frau des \dots\ or Schwester von \dots''.
\end{reviewercomment}
\response Added.
\where{p.~12, l.~437.}

\item
\begin{reviewercomment}
Line 427: provide a translation for ``Tochter, Frau des \dots''.
\end{reviewercomment}
\response Added.
\where{p.~14, l.~465.}

\item
\begin{reviewercomment}
Line 467: provide a translation for ``[\dots] e.g.\ (Frau des Heinrich), (Mutter der Antonie), (Mann
der Marianne), [\dots]''.
\end{reviewercomment}
\response Added.
\where{p.~16, ll.~505--506.}

\item
\begin{reviewercomment}
Line 476: provide a translation for ``[\dots] (Tochter des Ernst Ferdinand Klein) [\dots]''.
\end{reviewercomment}
\response Added.
\where{p.~16, ll.~515--516.}

\item
\begin{reviewercomment}
Table 6: provide translations for ``Frau des'', ``Baronin von'', ``Arzt in Berlin'' and ``wohnhaft in
der Friedrichsgracht 50''.
\end{reviewercomment}
\response Added, together with a translation of ``Fournier, L.\ und Compagnie''.
\where{Table~6, p.~17.}

\item
\begin{reviewercomment}
Lines 532 ff.: a discussion incorporating Section 3 (Case Study) and elements of Section 4
(Limitations) would be very helpful and is needed, especially from an RSE perspective -- for example,
how the NFDI and the basic service TS4NFDI come into play, and what that would change.
\end{reviewercomment}
\response We added a Discussion section (Section~5). It uses the editorial review of the case study to
assess the certainty tiers: editors changed an existing label in 5.8\,\% of the entries with the tier
\texttt{low} but in 26.3\,\% of those with the tier \texttt{contested}, while all entries with the tier
\texttt{high} and all but one with the tier \texttt{single} remained unreviewed. From an RSE perspective, it argues that the
classifiers are the replaceable part of the workflow, that the stable part is the layer of
normalization, abstention handling, and provenance, and that external resources such as models need
versioning. It then describes how TS4NFDI could come into play: the mappings from each method's raw
labels onto the shared label set could be published as reusable terminology mappings instead of
remaining in the source code.
\where{p.~20, ll.~576--593.}

\item
\begin{reviewercomment}
Footnotes 9/10: add ``(accessed on DATE)'', or cite them as bibliography entries for the sake of RSE
visibility.
\end{reviewercomment}
\response We replaced both repository footnotes with citations of a bibliography entry for the
software (see also Reviewer~C, C.1).
\where{p.~19, ll.~573--574; p.~21, l.~625.}

\item
\begin{reviewercomment}
Line 551: what is ``\S 4''?
\end{reviewercomment}
\response It referred to Section~4; we now write ``Section~4''.
\where{p.~20, l.~612.}

\item
\begin{reviewercomment}
Line 588: ``Various LLMs were used for grammar and spelling suggestions [\dots]'' -- which LLMs
exactly, and for what purpose?
\end{reviewercomment}
\response The declaration now names the model families and version ranges (Anthropic Claude Opus
4.5--5 and Sonnet, OpenAI GPT 5.1--5.5, Google Gemini Pro and Flash) and states their purposes:
spelling, grammar, and phrasing suggestions during manuscript preparation, and code generation within
agentic coding workflows during software development.
\where{p.~21, l.~649 -- p.~22, l.~653.}

\item
\begin{reviewercomment}
Footnotes 11--14: add ``(accessed on DATE)''.
\end{reviewercomment}
\response Added to the former footnotes 11--13, now footnotes 12--14. The former footnote 14 (the
repository link) was replaced by the bibliography entry for the software [KNW26], which includes the
access date.
\where{p.~21, footnotes 12--14.}

\item
\begin{reviewercomment}
Bibliography: use either a DOI or a link, not both [\dots]
\end{reviewercomment}
\response We removed the URL from all eight references that also had a DOI. References without a DOI
keep their URL.
\where{Bibliography, p.~22, l.~654 -- p.~24, l.~729.}

\end{enumerate}

% ======================================================================
\section*{Reviewer K}

Reviewer K's comments concern running the application and the missing register data; they are
answered under ``Software release'' above.

% ======================================================================
\section*{Further changes}

While addressing the comments, we made the following additional corrections:
\begin{itemize}
  \item \textbf{Table~6} was checked against the case-study output. We corrected the certainty tiers
        of S0001317, S0001502, and S9437075, showed the resulting labels of S0001502 and S9437075
        before review, corrected the vote counts for S0001932 and S0001317, and corrected the number
        in the caption (850 instead of 854 entries).
  \item \textbf{Certainty terminology:} the internal certainty tier and the TEI \texttt{@cert} value
        are now distinguished throughout, including the voting paragraph, Section~3.5.3, the LLM
        output schema, and the TEI example, whose \texttt{@cert} value is now \texttt{unknown}.
  \item \textbf{References:} the TEI Guidelines are cited as one entry (TEI P5, version 4.12.0) with
        page-specific locators; entries for the ediarum manual and the TS4NFDI architecture were added.
  \item \textbf{Section numbering:} the new Discussion is Section~5; the Conclusion is now Section~6.
\end{itemize}

\end{document}
